\section*{Abstract}
\label{sec:abstract}
\noindent\rule{\linewidth}{0.8pt}\par\nopagebreak
\vspace{2pt}
Interaktive virtuelle Prototypen ermöglichen ein frühzeitiges Testen des Verhaltens digitaler Objekte. Vorhandene 3D-Assets enthalten vorrangig Geometrie, Transformationen und Materialien, aber keine explizite Interaktionslogik. Diese Arbeit untersucht, wie sich solche statischen 3D-Objekte mithilfe natürlicher Sprache in interaktive Prototypen überführen lassen, ohne ihre Geometrie zu ersetzen. Dafür wurde der \textit{Vivian Specification Generator} (VSG) konzipiert und prototypisch umgesetzt. Der Ansatz verwendet eine Unity-basierte Vorverarbeitung sowie eine modulare, \textit{Large Language Model} (LLM)-gestützte Pipeline bestehend aus \textit{Scene Analyzer}, \textit{Interaction Planner}, \textit{Specification Generator}, \textit{Consistency Reviewer} und \textit{Unity-Validator}. Aus strukturierten Szenendaten, gerenderten Ansichten und Nutzer*innenbeschreibungen erzeugt der VSG Vivian-kompatible Spezifikationsdateien für ein zustandsbasiertes Interaktionsverhalten. Die technische Systemevaluation anhand der Testassets Lampe, Display und Toaster mit insgesamt 60 Durchläufen zeigt, dass der VSG vorhandene 3D-Objekte grundsätzlich in interaktive Prototypen überführen kann, da 57 Spezifikationen technisch valide generiert wurden. Die implementierten Korrekturmechanismen erhöhten die technische Erfolgsquote außerdem von 88,3 auf 95,0 Prozent.
Gleichzeitig verdeutlichen die Ergebnisse, dass technische Validität nicht mit funktionaler Korrektheit gleichzusetzen ist.
\par\nopagebreak\vspace{2pt}
\noindent\rule{\linewidth}{0.8pt}

\vspace{1em}
\noindent\textbf{Schlagwörter:} Interaktive Prototypen, KI-gestützte Transformation, Virtual Prototyping, Generative KI, Large Language Models, Virtual Reality
